\documentclass[11pt,letterpaper]{article}
\usepackage[T1]{fontenc}
\usepackage{lmodern}
\usepackage[margin=0.88in]{geometry}
\usepackage{amsmath,amssymb,amsthm,mathtools,mathrsfs}
\usepackage{microtype,booktabs,tabularx,array,longtable,enumitem}
\usepackage{xcolor,xurl,fancyhdr,needspace}
\usepackage[colorlinks=true,linkcolor=blue!45!black,citecolor=blue!45!black,urlcolor=blue!50!black]{hyperref}
\hypersetup{pdftitle={Refining the ETH contraction: a residue-aware Cauchy basis and a cheaper recursive schedule},pdfauthor={Research prepared for Vladimir Reshetnikov},pdfsubject={Explicit exponential certificates, exact local checks, and a review of seven ETH proposal packets}}
\setlength{\headheight}{14pt}
\pagestyle{fancy}
\fancyhf{}
\fancyhead[L]{\small Refining the ETH contraction}
\fancyhead[R]{\small 26 October 2026}
\fancyfoot[C]{\thepage}
\setlength{\parindent}{0pt}
\setlength{\parskip}{0.35em}
\setlength{\emergencystretch}{3em}
\setlist{nosep,leftmargin=1.6em}
\allowdisplaybreaks
\newtheorem{theorem}{Theorem}[section]
\newtheorem{lemma}[theorem]{Lemma}
\newtheorem{proposition}[theorem]{Proposition}
\newtheorem{corollary}[theorem]{Corollary}
\theoremstyle{definition}\newtheorem{definition}[theorem]{Definition}
\theoremstyle{remark}\newtheorem{remark}[theorem]{Remark}
\newcommand{\poly}{\operatorname{poly}}
\newcommand{\polylog}{\operatorname{polylog}}
\newcommand{\vTwo}{v_2}
\newcommand{\Z}{\mathbb Z}
\newcommand{\Q}{\mathbb Q}
\newcommand{\F}{\mathbb F}
\newcommand{\OO}{\mathcal O}
\newcommand{\LL}{\mathcal L}
\newcommand{\ceil}[1]{\left\lceil#1\right\rceil}
\newcommand{\floor}[1]{\left\lfloor#1\right\rfloor}
\newcommand{\ppin}{aa6933442eb936962de79ba6d321cd0391c450bd}
\newcommand{\opin}{7ef5c0cc5ca4b0e072b7a050fca09ab24bc6fa3a}
\newcommand{\newconstant}{133.20737249245\ldots}
\title{\vspace{-1.1em}Refining the ETH contraction\\[0.4em]
\large A residue-aware Cauchy basis, sharper packing,\\and a cheaper recursive precision schedule}
\author{Research prepared for Vladimir Reshetnikov}
\date{26 October 2026}
\begin{document}
\maketitle
\vspace{-0.8em}
\begin{abstract}
\noindent
We review all seven ETH proposal archives present in the pinned ProveIt
snapshot, together with the cited OpenAI preprint and the separately supplied
49-page manuscript. We obtain a new explicit contraction certificate
\[
 (s+\LL)^{O(1)}\,2^{(\newconstant+o(1))s/\sqrt{\log_2 s}}
\]
for the same normalized Boolean tensor interface used by the strongest
previous explicit schedule. Its displayed coefficient was $648.4410666487\ldots$;
combining that schedule with the prior factorial-Cauchy basis already gives
$374.3776243832\ldots$. Our proof improves both. The contributions are an
asymptotically sharp use of the packing load bound, an additive precision
guard replacing doubled packet exponents, and an exact Cauchy separation
that exploits a missing residue class in the actual rows. The last changes
the mode-weight slope from $3$ to $64/21$ without adding mode families.
The new local lemmas are proved here; the complete contraction conclusion
uses the unchanged deformation and recursive identities of the preprint.
This is an improvement of a specified exponential upper-bound coefficient.
It does not establish a better logarithmic order for unrestricted 3-SAT,
verify the entire preprint, or compare unknown constants in the separate
49-page algorithm. Exact executable checks accompany the proofs.
\end{abstract}

\section{The result and the comparison standard}
Let $n$ be the number of declared Boolean variables, $m$ the input clause
count, and $k$ a fixed clause width. The requested case is $k=3$. All
algorithms discussed here count satisfying assignments exactly and therefore
also decide satisfiability, subject to their stated hypotheses. Polynomial
factors include reading the encoded input. Unless explicitly written as
$\ln$, logarithms are base two.

The strongest statement in the cited OpenAI manuscript is its quantitative
Appendix D, incorporated into its main theorem:
\begin{align}
 T_{\mathrm{general}}(n,m)&=(n+m)^{O_k(1)}
  2^{O_k(n\log\log n/\sqrt{\log n})},\label{eq:general}\\
 T_{\mathrm{fixed}}(n,m)&=(n+m)^{O_{k,D}(1)}
  2^{O_{k,D}(n/\sqrt{\log n})}\qquad(m\le Dn).\label{eq:sparse}
\end{align}
Both use polynomial space~\cite{openai,reduction}. The earlier
$(\log\log n)^2$ numerator in the body and the reviewed Lean statement is
a weaker endpoint. Reproducing Appendix D would not improve the strongest
existing proposal.

Our principal theorem concerns the contraction subroutine with a normalized
wire budget $s$ and description length $\LL$; its full interface is in
Section~\ref{sec:interface}. Write $\lambda=\ceil{\log_2 s}$. The result is
\begin{equation}\label{eq:headline}
 T_{\mathrm{new}}(s,\LL)\le (s+\LL)^{O(1)}
 2^{(C_*+o(1))s/\sqrt\lambda},\qquad
 C_* = \frac{9\pi}{4\ln2}\sqrt{\frac{1365}{8}}
     =\newconstant.
\end{equation}
The $o(1)$ is uniform over instances satisfying fixed interface constants.
For example, $C_*+o(1)$ can be replaced by $134$ for all sufficiently large
$s$. The threshold is not asserted to be practical.

\begin{table}[ht]
\centering\small
\begin{tabularx}{\textwidth}{@{}X r@{}}
\toprule
Specified contraction certificate & Coefficient of $s/\sqrt{\log_2s}$\\
\midrule
Previous retuned schedule and improved taper~\cite{schedulepacket}
 & $648.4410666487$\\
That schedule composed with the existing factorial-Cauchy basis
 & $374.3776243832$\\
Boundary infimum of the earlier optimized family, with that basis
 & $368.6540181167$\\
New packing and additive guard, existing slope-$3$ basis
 & $134.2604115868$\\
New packing, additive guard, and the residue-aware basis
 & $133.2073724925$\\
\bottomrule
\end{tabularx}
\caption{Matched normalized interfaces and bit-cost accounting. These are
upper-bound coefficients, not measured times or lower bounds for earlier
algorithms. Rows two and three compose the earlier schedules with the prior
basis; they are not quoted standalone results.}\label{tab:constants}
\end{table}

The new certificate is smaller than every explicit comparable contraction
certificate established in the seven reviewed packets, including the
displayed earlier schedule composed with factorial-Cauchy modes.
For completeness, the earlier schedule packet also optimizes its
continuous parameter family: after composition with that basis its
boundary infimum is approximately $368.654018$, near $r=251$.
Strict packing slack approaches rather than attains that boundary.
The new coefficient is smaller than this comparison too. Its logarithmic
order is still $s/\sqrt{\log s}$. Consequently
\eqref{eq:general}--\eqref{eq:sparse} remain the established SAT orders.
Several packets propose faster orders under additional unproved compression
or summation assumptions. We do not claim to surpass those hypothetical
algorithms, to prove their missing assumptions, or to establish a runtime
inequality between implementations whose constants have not been specified.

\section{Repository review and provenance}\label{sec:review}
The ProveIt snapshot is
\begin{center}\small\texttt{\ppin}.\end{center}
The cited OpenAI source snapshot is
\begin{center}\small\texttt{\opin}.\end{center}
The three relevant intake commits below are dated 26 October 2026 in
Pacific time. Their contents, including the ZIP members rather than only
commit messages, were inspected. The archive manifest records immutable
URLs, Git object identifiers where available, byte sizes, and SHA-256
fingerprints. No repository files or remote branches were changed.

A subsequent ETH placement commit, \texttt{2b925d8800}, moved these
materials into \path{Computability/SubexponentialSAT/Research}.
That commit and the new project README were also reviewed. All 119
placed packet files were independently matched against the downloaded
archive members by their Git blob hashes; the mathematical corpus is
unchanged. The report retains the earlier immutable archive links, and
the provenance manifest records the current paths as well~\cite{placement}.

\begin{table}[ht]
\centering\small
\begin{tabularx}{\textwidth}{@{}p{0.19\textwidth}X@{}}
\toprule
Intake commit & ETH-related additions\\
\midrule
\texttt{a505bd768b} & Factorial-Cauchy modes; contraction/taper schedule;
 density-aware counting; original \texttt{paper.pdf}; separate
 \texttt{eth.pdf}.\\
\texttt{64b3202ee8} & Precision refinement; scalar parity/refinement packet;
 second density refinement.\\
\texttt{8afdbbea48} & Asymptotic-analysis packet with conditional
 degree-preserving compression and restricted method barriers.\\
\bottomrule
\end{tabularx}
\end{table}

\Needspace{6\baselineskip}
\subsection{What each prior packet actually establishes}
\begingroup\small
\begin{longtable}{@{}>{\raggedright\arraybackslash}p{0.22\textwidth}>{\raggedright\arraybackslash}p{0.43\textwidth}>{\raggedright\arraybackslash}p{0.29\textwidth}@{}}
\toprule
Packet & Established contribution within its stated model & Stronger unresolved target\\
\midrule\endfirsthead
\toprule Packet & Established contribution within its stated model & Stronger unresolved target\\\midrule\endhead
Factorial Cauchy~\cite{cauchypacket} & Integral exact separation with
 weights $1+b+2v_2(b!)=3b+O(\log b)$; matched entropy coefficient reduced
 by $1/\sqrt3$. & Better log-power needs additional structure.\\[0.5em]
Contraction improvements~\cite{schedulepacket} & Exact Hermite condition
 count, leading precision $6MD/u$, schedule coefficient $648.44\ldots$,
 and streamed packet coefficient extraction. & Same global log-power.\\[0.5em]
SAT density improvement~\cite{densitypacket} & For fixed-polylogarithmic
 density $B$, exponent $O(n\log(2+B)/\sqrt{\log n})$, with explicit
 row-construction costs. & Uniform compression at arbitrary density.\\[0.5em]
ETH density refinement~\cite{densitypacket2} & A related density-sensitive
 construction and checks of retuning and interpolation losses. & No proved
 unrestricted lower log-order.\\[0.5em]
Precision refinement~\cite{precisionpacket} & Exact node-dependent
 precision and optimal taper analysis for its loss functional. & A
 structured-mode summation hypothesis would give an $s/\log s$ scale.\\[0.5em]
Scalar parity refinement~\cite{scalarpacket} & Removes the nilpotent
 variable in that reduction and lowers incidence degree from
 $k(a+2n)$ to $k(a+n)$. & A second contraction needs a compact
 high-degree-group factorization.\\[0.5em]
Asymptotic analysis~\cite{analysispacket} & Restricted counting and
 encoding barriers; precise interface hypotheses. & A degree-preserving
 compiler would yield $O(n\sqrt{\log\log n/\log n})$;
 a stronger compiler would yield $O(n/\sqrt{\log n})$.\\
\bottomrule
\end{longtable}
\endgroup

These results are not seven independent proofs of a new log-power. In
particular, an exact finite representation of a mode sum does not supply
an efficient contraction of that representation. Likewise a retained
incidence-degree cap does not automatically bound the multilinear degree
of a binary subset index. The reviewed packets generally state these
boundaries explicitly; this report preserves them.

\subsection{The separate 49-page manuscript}
The file \texttt{eth.pdf}, titled \emph{A $2^{o(n)}$ time algorithm for
3-SAT}, attributes itself to Claude/Anthropic and is dated 26 October 2026,
3:13 PM PT. It is a different construction, based on decompositions of a
crossing, recursive clique algorithms, and partition dimensions. Its
Theorem 1 has exponent
\[
 O\!\left(\min\left\{
  k^2n\frac{\log\log n}{\sqrt{\log n}},\,
  km\frac1{\sqrt{\log m}}
 \right\}\right)
\]
in its stated computational model. For fixed $k$ this has the same leading
orders as the OpenAI manuscript. Its Appendix B supplies
$\rho=1/\sqrt2$ to the recurrence in Appendix A, which yields the
logarithmic exponent $\log_2(1/\rho)=1/2$~\cite{otherpaper}.
It does not provide a coefficient directly comparable with
Table~\ref{tab:constants}'s normalized Boolean-wire interface. We therefore
record its rate and distinct model without claiming that
\eqref{eq:headline} is faster than its complete algorithm. Neither that
manuscript's Lean development nor the OpenAI Lean development was rebuilt
for this report.

\section{The contraction interface and inherited machinery}\label{sec:interface}
The normalized input is a bipartite Boolean tensor multigraph with at most
$s$ edges, polynomially many vertices in $s$, and maximum degree at most
$\Delta=\ceil{s^{9/10}}$. Parallel edges have separate ports. The entries
belong to
\[
 \mathcal B_{P_0,\ell}=(\Z/2^{P_0}\Z)[V]/(V^\ell),
 \qquad 1\le P_0,\ell\le s.
\]
Given its immutable description and a Boolean assignment to its ports,
an original tensor entry is computable in
$(s+\LL)^b2^{A s^{9/10}}$ bit operations and $(s+\LL)^b$ bits of
space for fixed $A,b$. Constants in a more general $Cs$ interface can be
absorbed into $s$; the explicit coefficient then refers to that rescaled
$s$, not directly to $n$.

The source extends scalars to an unramified extension
$\OO$ of $\Z_{(2)}$, with sufficiently many distinct residue-field
elements for unit-difference interpolation. Valuations are normalized by
$v_2(2)=1$ and are taken coordinatewise in the truncated $V$ ring.
All precisions and extension dimensions used below are polynomial in $s$.

At a recursive call with $M$ wires, a packing creates packets of
$L=2^a$ slots. A Cauchy polynomial matrix $W(T)$ has $W(0)=I$ and
degree $2L-1$ per entry. Applying the multilinear row tensors to
$W(T)X$ gives a deformation $F(T)$ of degree at most
$M(2L-1)$ whose value at zero is the desired contraction. At
\[
 T=t+\alpha\zeta,\quad
 \alpha=2^{a+1},\quad t\in\{-L,\ldots,L-1\},\quad\zeta\in\OO,
\]
the local kernel is integral. For retained $i=v_2(t)$, grouping vertices
by label coordinate $i$ makes $W(t)$ local. The remaining crossing
interactions have an integral separated expansion into modes.

The source's recursive identity sums over mode exponent tuples and takes
a product of one smaller child contraction per group. Its validity needs
only integral target coefficients with known valuations and source
features confined to a single group. These are exactly the properties
preserved by our new basis. We retain the deformation, recovery identity,
and child construction~\cite{contraction}.

\paragraph{Immutable precision semantics.}
A tensor is tagged with its creation precision. Whenever a later call
uses it at a larger precision, its entry is first evaluated at its creation
precision and the same integer coordinate representatives are lifted.
Canonical lifting is not asserted to preserve multiplication or addition.
New tensors retain their creation-time descriptions. This rule is
necessary for the source's induction and is unchanged here.

\section{The taper improvement already available}\label{sec:taper}
We use the strongest established precision refinement rather than counting
it as a new contribution. Put $q_i=h(u-i)$ for $0\le i<u$, and $q_i=0$
otherwise, where $1\le u\le a$. There are $L/2^i$ nonzero grid nodes of
valuation $i<u$. Consequently the exact number of Hermite conditions is
\begin{equation}\label{eq:hermitecount}
 \sum_t q_{v_2(t)}=hL(2u-2+2^{1-u}).
\end{equation}
Use the integer choice
\begin{equation}\label{eq:hchoice}
 h=\ceil{\frac{M(2L-1)+1}{L(2u-2+2^{1-u})}}.
\end{equation}
It supplies more conditions than the deformation degree. Whenever
$u\longrightarrow\infty$, $u/M\longrightarrow0$ and $L\longrightarrow\infty$
uniformly across active calls,
\begin{equation}\label{eq:hscale}h=(1+o(1))M/u.\end{equation}

For $D_i=L/2^i$, the exact loss used in the precision packet is
\[
 B_i=D_i\sum_{j=1}^{a-i}\frac{j}{2^j}(q_i-q_{i+j}).
\]
The sufficient working precision is
\begin{equation}\label{eq:exactprecision}
 P_t=p+B_i+(q_i-1)(a+1-i),
\end{equation}
where $p$ is the call's output precision. The Hermite multiplier has
valuation $-B_i$, and evaluation of a coefficient of order $j<q_i$
at $-t/\alpha$ loses $j(a+1-i)$ powers of two. These statements prove
\eqref{eq:exactprecision} by bounding the contribution of each input error
separately. The truncated reciprocal with constant term one is integral,
so forming it loses no precision~\cite{precisionpacket}.

Since $q_i-q_{i+j}\le hj$ and
$\sum_{j\ge1}j^2/2^j=6$, we have $B_i\le6hD_i$. Once the inherited
precision and the final jet-evaluation term are negligible, this gives
\begin{equation}\label{eq:precisionsix}
 P_t\le(6+o(1))MD_i/u.
\end{equation}
The next sections establish that negligibility for the new schedule.
The exact optimal-taper result in the earlier precision packet explains
why changing only this taper does not supply another log-order gain.

\section{New contribution I: keep the sharp packing bound}\label{sec:packing}
Let $r$ be a fixed integer, let each vertex label belong to $[r]^a$, and
put $N=r^{2a}$, the number of ordered pairs of labels. The source's
conditional-expectation argument produces label-class loads satisfying
\[
 \max_\gamma X_\gamma\le M/N+\sqrt{2M\Delta}.
\]
Padding every label-pair class to the maximum load rounded up to a multiple
of $L$ therefore gives
\begin{equation}\label{eq:sharp-pad}
 H_{\rm pad}\le M/N+\sqrt{2M\Delta}+L.
\end{equation}
Define
\[
 \epsilon_s(M,a)=\sqrt{\frac{2N^2\Delta}{M}}+\frac{LN}{M}.
\]
Multiplication of \eqref{eq:sharp-pad} by $N$ yields
$H_{\rm pad}N\le(1+\epsilon_s)M$. Thus the same deterministic
packing has the following sharper consequences:
\begin{align}
 \#\text{packets}&\le(1+\epsilon_s)M/L,\label{eq:packetsharp}\\
 \#\text{ports per side in a group}&\le(1+\epsilon_s)M/r,\\
 M_{\rm child}&\le2(1+\epsilon_s)M/r,\label{eq:childsharp}\\
 c_t:=\#\text{column-tile families}&\le2(1+\epsilon_s)M/D_i.
 \label{eq:colsharp}
\end{align}
The last line uses $2^{i+1}$ possible tiles per packet. It counts mode
families, each of which has one variable per mode index. In contrast,
the source and earlier schedules replace the right side of
\eqref{eq:sharp-pad} by $3M/N$, giving $6M/D_i$ families and
$6M/r$ child wires. Under strict-power packing slack,
$\epsilon_s=o(1)$ uniformly, so that factor three can be removed from
the leading constant. This is a refinement of the proven load estimate;
the packing algorithm itself does not change.

\section{New contribution II: an additive precision guard}\label{sec:schedule}
Take
\begin{equation}\label{eq:schedule}
 \begin{gathered}
 \lambda=\ceil{\log_2s},\quad r=54,\quad
 H=\ceil{s/\lambda},\quad
 g=\bigl(\ceil{\log_2\lambda}\bigr)^2,\\
 u_d=\floor{\frac{\lambda}{260\cdot9^d}},\qquad
 a_d=\sum_{j=0}^d u_j+(d+1)g,
 \qquad a_{-1}=0.
 \end{gathered}
\end{equation}
An active call has $M>H$; smaller calls are evaluated directly.
As in the source, bounded exceptional input sizes are handled by exhaustive
summation. The relevant integer inequalities can also be checked before
recursion. The following proof shows that the recursive regime applies
uniformly once $s$ exceeds a fixed threshold.

\subsection{Packing slack and a noncircular induction}
Write $A_0=9/2080$. At depths $O(\log\lambda)$,
\begin{equation}\label{eq:a-budget}
 a_d\le A_0\lambda+O((\log\lambda)^3).
\end{equation}
The exact integer inequality $54^{25}<2^{144}$ gives
\begin{equation}\label{eq:powergap}
 \frac9{10}+4A_0\log_2 54
 <\frac9{10}+4\frac9{2080}\frac{144}{25}
 =\frac{3249}{3250}<1.
\end{equation}
Therefore $N^2\Delta=s^{1-\Omega(1)}$ and $LN=s^{1-\Omega(1)}$,
uniformly at those depths. Since $M>s/\lambda$, we get
$\epsilon_s=s^{-\Omega(1)}$. Also $2L\le\Delta$ eventually.

To justify the depth range without assuming the desired shrinkage, start
with the provisional invariant $M_d\le s/16^d$. It holds at the root.
At an active call it bounds $d$ by $\log_{16}\lambda+O(1)$, so
\eqref{eq:a-budget} and the packing estimates already apply. Because
$2/54<1/16$, \eqref{eq:childsharp} implies the invariant at the next
depth for all sufficiently large $s$. This closes the induction.
Since $9<16$, this provisional argument also proves that the retained
$u_d$ tend to infinity uniformly.

We may now use the sharp ratio. Multiplying
$M_{d+1}\le(1+\epsilon_s)M_d/27$ over $O(\log\lambda)$ depths,
with a uniform negative-power bound on $\epsilon_s$, gives
\begin{equation}\label{eq:sharppath}
 M_d\le(1+o(1))s/27^d.
\end{equation}
At active calls this implies
$d\le\log_{27}\lambda+O(1)$ and
\begin{equation}\label{eq:ulower}
 u_d=\Omega(\lambda^{1/3}),\qquad u_d/M_d=o(1).
\end{equation}
Floors in \eqref{eq:schedule} hence cause only uniform relative $o(1)$
errors. Both \eqref{eq:hscale} and the sharp-packing estimates hold
throughout the recursion.

\subsection{The guard absorbs inherited precision}
The critical identity is exact:
\[
 a_d-u_d=a_{d-1}+g.
\]
For every retained valuation $i<u_d$,
\begin{equation}\label{eq:Dlower}
 D_i=2^{a_d-i}\ge2^{a_{d-1}+g+1}.
\end{equation}
The source precision argument, with the smaller increments from
\eqref{eq:exactprecision}, gives
\[
 p\le s\lambda^{O(1)}2^{a_{d-1}}.
\]
This follows directly by summing the increments along the path; each
earlier increment is at most $s\lambda^{O(1)}2^{a_{d-1}}$, and the
path has $O(\log\lambda)$ levels. Hence, at an active call,
\[
 \frac{p}{MD_i/u_d}\le\lambda^{O(1)}2^{-g}=o(1).
\]
Also
\[
 \frac{hu_d(a_d+1)}{MD_i/u_d}
 =O\!\left(\frac{u_d(a_d+1)}{D_i}\right)
 \le\lambda^{O(1)}2^{-g}=o(1).
\]
Thus \eqref{eq:precisionsix} follows uniformly. The additive guard
contributes only $O((\log\lambda)^3)=o(\lambda)$ to the packet
exponent. It replaces the earlier choice $a_d=2\sum_{j\le d}u_j$,
whose extra full copy of the retained-level budget is unnecessary.

\subsection{Sizes and entry costs}
All packet and label spaces remain polynomial in $s$ by
\eqref{eq:a-budget}. The deformation degree is
$O(sL)=s^{1+A_0+o(1)}<s^4$, so the source's residue field of
size $2^{4\lambda}$ still supplies enough interpolation points.
All precisions are polynomial in $s$. New packet degrees are at most
$2L\le\Delta$, and old vertex degrees are unchanged. Even the source's
original packet-entry procedure costs only
$\exp(O(L\log s))=\exp(o(s^{9/10}))$ here. The previously proved
streamed finite-difference routine can be retained, but is not needed to
justify our leading coefficient.

\section{A residue-aware Cauchy separation}
\label{sec:pivot21}

This section changes the separated basis used for crossing entries in one
contraction call. It does not assume an unproved faster summation routine.
The new basis has the same source and target groups as the previous
factorial-Cauchy basis, but its guaranteed valuation grows with slope
$64/21$ instead of $3$.

\subsection{The actual row domain}

Let $\mathcal O$ be the unramified coefficient ring used in the source
construction, with valuation $\nu_2(2)=1$. In particular, its reduction
modulo two can be a proper extension of $\mathbb F_2$. Set
\[
 L=2^a,\quad i=\nu_2(t)<a,\quad
 \alpha=2^{a+1},\quad T=t+\alpha\zeta,\quad \zeta\in\mathcal O,
 \qquad J=2^{a-i-1}.
\]
A column tile has positions
$y=y_0+2^{i+1}j$, where $0\le y_0<2^{i+1}$ and $0\le j<J$.
For a crossing row $x$, the bits in position $i$ of $x$ and $y_0$
are equal. Define
\[
 d_0=t-y_0+x,\qquad d=d_0+\alpha\zeta,\qquad
 z_x=\frac{2^{i+1}}{d},\qquad w_x=\frac{z_x}{2},
 \qquad R=a+1-i.
\]
The source's crossing condition gives $d_0\ne0$ and
$\nu_2(d_0)\le i$. With
\[
 \mathscr C(T)=\prod_{q=-L}^{L-1}(T-q),\qquad
 c_x=\frac{\mathscr C(T)}{\mathscr C'(0)d},
\]
the kernel is
\begin{equation}
 W_{xy}(T)=\frac{c_x}{1-z_xj},
 \qquad z_x\in2\mathcal O,\quad \nu_2(c_x)\ge1.
 \label{eq:pivot21-kernel}
\end{equation}
The last lower bound may be sharpened to $R$, but an offset of one
already suffices for the comparison below.

\begin{lemma}[A missing residue class]
\label{lem:pivot21-domain}
Put $w_0=2^i/d_0\in\mathbb Z_{(2)}$. Then
\[
 w_x-w_0\in2^R\mathcal O.
\]
If $i\ge1$, define $c\in\{2,6\}$ by
\[
 c=\begin{cases}
       6,&\text{bit }i-1\text{ of }y_0\text{ is zero},\\
       2,&\text{bit }i-1\text{ of }y_0\text{ is one}.
    \end{cases}
\]
For $R\ge3$, the residue of $w_x$ modulo eight lies in the prime
subring and is different from $8-c$. If $i=0$, $w_x$ is odd and
either choice of $c$ is valid.
\end{lemma}

\begin{proof}
Write $v=\nu_2(d_0)\le i$. Adding $2^{a+1}\zeta$ preserves that
valuation. Therefore
\[
 w_x-w_0=
 -\frac{2^{i+a+1}\zeta}{d_0(d_0+2^{a+1}\zeta)},\qquad
 \nu_2(w_x-w_0)\ge i+a+1-2v\ge R.
\]
This proves the prime-subring assertion when $R\ge3$.

Only the case $\nu_2(w_0)=1$ could give residue two or six. In this
case $v=i-1$. Write the lower $i$ bits of $x$ and $y_0$ as $x_<$
and $y_<$. Since their bits in position $i$ agree, the other part of
$x-y_0$ is divisible by $2^{i+1}$. The difference $x_<-y_<$ lies
strictly between $-2^i$ and $2^i$ and, by $v=i-1$, is an odd
multiple of $2^{i-1}$. It must be $+2^{i-1}$ when bit $i-1$ of
$y_0$ is zero, and $-2^{i-1}$ otherwise. Since
$t/2^{i-1}\equiv2\pmod4$, the odd number
$d_0/2^{i-1}$ is respectively three or one modulo four. Its inverse
has the same residue, so $w_0$ is respectively six or two modulo
eight. In all other valuation cases neither residue can occur.
For $i=0$, $d_0$ is odd, which proves the last statement.
\end{proof}

The exceptional value $R=2$ gives $J=1$. It has only the mode $b=0$
and needs no residue exclusion.

\subsection{An explicit sequence of auxiliary row points}

Repeat the word of length 21
\[
 A\ B\ A\ B\ C\ A\ B\ A\ B\ C\ A\ B\ A\ A\ C\ A\ B\ A\ A\ C\ A.
\]
At the $k$th occurrence of a letter, starting at $k=0$, set
\begin{equation}
 v_h=\begin{cases}
       1+2k,&\text{letter }A,\\
       4k,&\text{letter }B,\\
       c+8k,&\text{letter }C,
     \end{cases}
 \qquad \rho_h=2v_h.
 \label{eq:pivot21-points}
\end{equation}
The index $h$ counts positions in the repeated word. Let $A_b,B_b,C_b$
be the numbers of the three letters among its first $b$ positions.
Thus a complete block has counts $(11,6,4)$. Define
\begin{equation}
 \eta_b=\min\left\{
 A_b+\nu_2(A_b!),\quad
 2B_b+\nu_2(B_b!)+C_b,\quad
 3C_b+\nu_2(C_b!)+B_b
 \right\}.
 \label{eq:pivot21-eta}
\end{equation}

\begin{lemma}[Factorials near the prime subring]
\label{lem:pivot21-near-prime}
If $q\in\mathcal O$, $q_0\in\mathbb Z_{(2)}$, and
$q-q_0\in2^Q\mathcal O$, then for every integer $0\le m<2^Q$,
\[
 \nu_2\!\left(\prod_{k=0}^{m-1}(q-k)\right)\ge\nu_2(m!).
\]
\end{lemma}

\begin{proof}
For $1\le r\le\lfloor\log_2m\rfloor$, the rational $q_0$ has a
well-defined integer residue modulo $2^r$. Among $m$ consecutive
integers, at least $\lfloor m/2^r\rfloor$ have that residue.
Since $r\le Q$, the corresponding factors $q-k$ are divisible by
$2^r$. Summing over $r$ proves the result by Legendre's formula.
For $m=0,1$ the assertion is immediate.
\end{proof}

This lemma is necessary: an arbitrary extension element of $\mathcal O$
need not have integral binomial coefficients. The actual rows have the
required proximity to the prime subring.

\begin{lemma}[Valuation of the pivot product]
\label{lem:pivot21-product}
For every actual crossing row and every $0\le b<J$,
\[
 \nu_2\!\left(\prod_{h=0}^{b-1}(w_x-v_h)\right)\ge\eta_b.
\]
Furthermore, $\eta_b$ is nondecreasing and
\[
 \eta_b=\frac{22}{21}b+O(\log(b+2)).
\]
\end{lemma}

\begin{proof}
Suppose first that $R\ge3$. The permitted residues split into three
cases. If $w_x$ is odd, the $A$ factors give
\[
 2^{A_b}\prod_{k<A_b}\left(\frac{w_x-1}{2}-k\right),
\]
while every $B$ or $C$ factor is a unit. Their valuation is at least
$A_b+\nu_2(A_b!)$.

If $w_x\equiv0\pmod4$, the $B$ factors similarly contribute at
least $2B_b+\nu_2(B_b!)$, each $C$ factor contributes exactly one
power of two, and the $A$ factors are units. The resulting bound is
the second term in \eqref{eq:pivot21-eta}.

The remaining permitted case is $w_x\equiv c\pmod8$. Now the $C$
factors contribute at least $3C_b+\nu_2(C_b!)$, each $B$ factor
contributes one power of two, and the $A$ factors are units. This
gives the third term.

To justify the factorial steps over $\mathcal O$, apply
Lemma~\ref{lem:pivot21-near-prime} after dividing by two, four, or
eight. Lemma~\ref{lem:pivot21-domain} supplies precisions $R-1$,
$R-2$, and $R-3$ respectively. The first two factor counts satisfy
\[
 A_b\le b<J=2^{R-2}<2^{R-1},\qquad
 B_b\le b<2^{R-2}.
\]
If $b=21q+r$, $0\le r\le20$, the positions of the letter $C$ give
\[
 C_b=4q+\lfloor r/5\rfloor\le\lfloor b/5\rfloor
 \le\lfloor b/4\rfloor.
\]
Thus, whenever $C_b>0$, $C_b<2^{R-3}$. If $R\le3$, no $C$ factor
occurs. The case $R=2$ has $b=0$, so its product bound is immediate.

Each expression in the minimum \eqref{eq:pivot21-eta} is
nondecreasing as the three counts increase. For $b=21q$,
Legendre's formula gives the exact identity
\[
 \eta_{21q}=22q-
 \max\{s_2(11q),s_2(6q),s_2(4q)\},
\]
where $s_2$ is the number of ones in the binary expansion.
For general $b$, the three counts differ by $O(1)$ from
$(11b/21,6b/21,4b/21)$. The bound
$s_2(m)=O(\log(m+2))$ proves the claimed asymptotic formula.
\end{proof}

\subsection{A finite separation with no nonunit division}

For arbitrary even auxiliary points $\rho_h$, set
\[
 U_b(z)=\prod_{h<b}\frac{z-\rho_h}{1-hz},\qquad
 V_b(j)=\prod_{h<b}\frac{j-h}{1-\rho_hj}.
\]
All denominators are units for $z\in2\mathcal O$ and integer $j$.

\begin{lemma}[Finite Cauchy elimination]
\label{lem:pivot21-finite}
For every nonnegative integer $q$,
\begin{equation}
 \frac1{1-zj}
 =\sum_{b<q}
 \frac{(1-b\rho_b)U_b(z)V_b(j)}{(1-bz)(1-\rho_bj)}
 +\frac{U_q(z)V_q(j)}{1-zj}.
 \label{eq:pivot21-finite}
\end{equation}
If $0\le j<J$ and $q=J$, the remainder is zero.
\end{lemma}

\begin{proof}
The elementary identity
\[
 (1-bz)(1-\rho_bj)-(z-\rho_b)(j-b)
   =(1-b\rho_b)(1-zj)
\]
shows that the $b$th summand is the difference of the consecutive
remainders $U_b(z)V_b(j)/(1-zj)$ and
$U_{b+1}(z)V_{b+1}(j)/(1-zj)$. Telescope. For $j<J$, the product
$V_J(j)$ contains the zero factor $j-j$.
\end{proof}

For the specific points \eqref{eq:pivot21-points}, define
\begin{align}
 \ell_b(j)&=
 \frac{(1-b\rho_b)\binom jb}
      {\prod_{h=0}^{b}(1-\rho_hj)},
 \label{eq:pivot21-source}\\
 r_b(x)&=
 c_x\,b!\,
 \frac{\prod_{h=0}^{b-1}(z_x-\rho_h)}
      {\prod_{h=0}^{b}(1-hz_x)}.
 \label{eq:pivot21-target}
\end{align}
Here $\binom jb=0$ for $b>j$.

\begin{theorem}[Residue-aware separated kernel]
\label{thm:pivot21-kernel}
For every actual crossing entry,
\[
 W_{xy}(T)=\sum_{b=0}^{J-1}r_b(x)\ell_b(j).
\]
Every source feature $\ell_b(j)$ belongs to $\mathbb Z_{(2)}$.
The target coefficient satisfies
\begin{equation}
 \nu_2(r_b(x))\ge\omega_b,
 \qquad
 \omega_b=1+b+\nu_2(b!)+\eta_b
          =\frac{64}{21}b+O(\log(b+2)).
 \label{eq:pivot21-weight}
\end{equation}
The weights $\omega_b$ are strictly increasing. The stronger offset
$R$ may replace one in this formula.
\end{theorem}

\begin{proof}
Use Lemma~\ref{lem:pivot21-finite} with $q=J$, and replace
$\prod_{h<b}(j-h)$ by the integer product $b!\binom jb$.
The integer binomial and the odd denominators make
\eqref{eq:pivot21-source} integral. In
\eqref{eq:pivot21-target}, factoring two from each of the $b$
numerator factors gives $2^b\prod_{h<b}(w_x-v_h)$. Combine
Lemma~\ref{lem:pivot21-product} with the explicit $b!$ and
$\nu_2(c_x)\ge1$. Legendre's formula gives the slope
$1+1+22/21=64/21$. Finally, $\eta_b$ is nondecreasing and the term
$b$ increases strictly, so the weights do also.
\end{proof}

The formula multiplies the factorial into the target. It never divides
by a factorial in the coefficient ring. Implementation uses integer
binomial coefficients, ring multiplication, and unit inversion only.
The points $\rho_h$ are auxiliary interpolation points; they need not
be actual rows in a packet.

\subsection{Consequences for the contraction accounting}

For each existing column tile, replace its source moment of degree $b$
by
\[
 L_b(X)=\sum_{\substack{\text{real columns }y\text{ in the tile}\\
                        y=y_0+2^{i+1}j}}
               \ell_b(j)X_y.
\]
There is still one mode of each weight per old tile. It is a linear
form in bits of the same source group and targets the same other group.
After expanding the row factors, a mode-exponent tuple $\beta$ whose
total certified weight is at least the working precision $p$ has
coefficient zero modulo $2^p$. Thus the retained condition is
\[
 \sum_{\text{tile families }f}\sum_{b=0}^{J-1}
             \omega_b\beta_{f,b}<p.
\]
The source's child factorization, wire count, degree bounds, and
creation-precision convention are preserved. Products defining the
features contain at most $J\le L$ factors of polynomial description
length, so the entry-oracle bounds remain valid. This argument also
permits the strongest compatible previous packet-entry routine.

To quantify the gain, let $N_\gamma(c,p)$ denote the abstract tuple
count with $c$ families and weights
$\omega_b=\gamma b+O(\log(b+2))$. Its standard generating-product
bound is
\[
 N_\gamma(c,p)\le
 \exp\!\left(\theta p+
       c\sum_{b\ge0}-\ln(1-e^{-\theta\omega_b})\right).
\]
As $\theta\downarrow0$, the sum has leading term
$\pi^2/(6\gamma\theta)$. Optimizing $\theta$ gives the leading
logarithmic bound
\[
 \ln N_\gamma(c,p)\le
 \left(\pi\sqrt{\frac{2}{3\gamma}}+o(1)\right)\sqrt{cp}.
\]
The actual finite cutoff only reduces this bound. In the active-call
regime, $p/c\to\infty$, so the displayed asymptotic applies uniformly
with the previously verified schedule.

Consequently, replacing the previous factorial-Cauchy weights of slope
$3$ by \eqref{eq:pivot21-weight}, while retaining a matched precision
and recursion schedule, multiplies the leading enumeration coefficient by
\begin{equation}
 \sqrt{\frac{3}{64/21}}=\sqrt{\frac{63}{64}}
   =0.9921567416\ldots.
 \label{eq:pivot21-ratio}
\end{equation}
Relative to the geometric weights of slope one, the factor is
$\sqrt{21/64}$. These are comparisons of explicit upper-bound
certificates. They are not wall-clock speedup measurements or lower
bounds on the old algorithm, and they do not change the logarithmic
order of the global SAT bound.

The new weights need not dominate the old weights at every finite
index. For example, the new weight at $b=4$ is $10$, while the old
factorial-Cauchy weight is $11$. The improvement is in the proved
leading asymptotic coefficient. If desired, an implementation may
count the old and new abstract tuple sets by ordinary generating-function
dynamic programming and choose the smaller set at each call. The numbers
of families, modes and precision units are polynomial in $s$, and the
counts have polynomial bit length, so this selection has polynomial
bit cost and space. It preserves both exact identities and the new
asymptotic upper bound. No uniform finite-input speedup is asserted.


\section{Entropy and the complete contraction theorem}\label{sec:global}
Suppose the common mode weights have the form
\[
 w_b=\gamma b+O(\log(b+2)),\qquad w_b\ge1,
\]
and are strictly increasing. The new kernel in the preceding section
gives $\gamma=64/21$. The old factorial-Cauchy basis gives $\gamma=3$.
There are at most $c_t=(2+o(1))M/D_i$ families. Extending each finite
family to infinitely many formal weights only enlarges the count.

For $\theta\downarrow0$, comparison with the integral
$\int_0^\infty-\ln(1-e^{-x})\,dx=\pi^2/6$ gives
\begin{equation}\label{eq:partition-asymptotic}
 \sum_{b\ge0}-\ln(1-e^{-\theta w_b})
 =\frac{\pi^2}{6\gamma\theta}+o(\theta^{-1}).
\end{equation}
One way to justify the perturbation is to fix any $\eta>0$: for all
sufficiently large $b$, $(\gamma-\eta)b\le w_b\le(\gamma+\eta)b$.
The finitely many omitted indices contribute $O(\ln(1/\theta))$.
The integral comparison squeezes $\theta$ times the sum between
$\pi^2/(6(\gamma+\eta))$ and $\pi^2/(6(\gamma-\eta))$; let
$\eta\downarrow0$. This proves the precise leading term without
requiring an unproved partition formula.

If $\mathcal N_t$ counts tuples satisfying
$\sum_{f,b}w_b\beta_{f,b}<P_t$, the generating product yields
\[
 \ln\mathcal N_t\le\theta P_t+
 c_t\sum_{b\ge0}-\ln(1-e^{-\theta w_b}).
\]
Use $P_t\le(6+o(1))MD_i/u_d$ and
$c_t\le(2+o(1))M/D_i$, and optimize the two leading terms.
The saddle scale is a constant times $\sqrt{u_d}/D_i$, which tends
uniformly to zero by \eqref{eq:Dlower}. We obtain
\begin{equation}\label{eq:localentropy}
 \ln\mathcal N_t\le
 \left(2\pi\sqrt{\frac2\gamma}+o(1)\right)
 \frac{M}{\sqrt{u_d}}.
\end{equation}
Integer ceilings on family counts are negligible because $M/D_i$
tends uniformly to infinity. The saddle estimate and the preceding
perturbation argument are therefore uniform across the recursion.

\begin{theorem}[Explicit refined contraction]\label{thm:main}
Assume the normalized interface of Section~\ref{sec:interface} and the
unchanged integral deformation and child identities of the source.
Use the exact taper \eqref{eq:hchoice}, precision
\eqref{eq:exactprecision}, schedule \eqref{eq:schedule}, sharp packing,
and the proved integral mode basis of slope $\gamma=64/21$.
Then deterministic exact contraction uses polynomial space and at most
\[
 (s+\LL)^{O(1)}2^{(C_\gamma+o(1))s/\sqrt\lambda}
 \quad\text{bit operations},\qquad
 C_\gamma=\frac{9\pi}{4\ln2}\sqrt{\frac{520}{\gamma}}.
\]
In particular $C_{64/21}=\newconstant$. The same new packing and
schedule with the existing slope-$3$ basis gives
$C_3=134.2604115868\ldots$.
\end{theorem}
\begin{proof}
Correctness follows by the same induction from leaf contractions as in
the source. The new mode identity is exact, all omitted coefficients
vanish at the working precision, and the source/target factorization
is unchanged. Hermite recovery has enough jets by
\eqref{eq:hermitecount}, at sufficient precision by
\eqref{eq:exactprecision}. Creation-time lifts remain fixed.

For cost, \eqref{eq:sharppath} and
$u_d=(1+o(1))\lambda/(260\cdot9^d)$ give
\[
 \sum_{d\ \mathrm{active}}\frac{M_d}{\sqrt{u_d}}
 \le(1+o(1))\sqrt{260}\frac{s}{\sqrt\lambda}
     \sum_{d\ge0}(3/27)^d
 =(1+o(1))\frac98\sqrt{260}\frac{s}{\sqrt\lambda}.
\]
Combining this with \eqref{eq:localentropy}, and dividing the natural
logarithm by $\ln2$, gives the asserted coefficient. Node choices,
interpolation values and group choices give only polynomial branching
per call, so their accumulated logarithm is
$O(\lambda\log\lambda)=o(s/\sqrt\lambda)$. A depth-first leaf sums
over at most $H=O(s/\lambda)$ bits. Its direct enumeration and entry
evaluation have logarithmic costs
$O(s/\lambda)+O(s^{9/10})=o(s/\sqrt\lambda)$.

Arithmetic, packing and local descriptions have polynomial cost per
generated term outside the permitted entry work. Because the branching
bound is independent of $\LL$, dependence on description length stays
a fixed polynomial rather than a depth-dependent polynomial power.
The finite Cauchy features require only polynomially many exact ring
operations and preserve the same packet-entry bound.

For space, enumerate one tuple and one interpolation value at a time,
accumulate a product of child values and then the running sum, and keep
only the current depth-first stack. There are $O(\log\lambda)$ levels.
Each holds polynomially many polynomial-bit data, immutable descriptions,
and residues. Exact rational recovery is the source's polynomial-bit
calculation. Thus space remains polynomial in $s+\LL$.
\end{proof}

\subsection{Why these schedule constants were chosen}
More generally, retain the old variance bound, use maximum degree
$s^\delta$ for fixed $\delta<1$, and set
$u_d=\floor{\varepsilon\lambda/b^d}$ with the same additive guard.
Provided $1<b<r/2$ and
\[
 \varepsilon<\frac{(1-\delta)(b-1)}{4b\log_2r},
\]
the leading coefficient is
\begin{equation}\label{eq:parameter-family}
 \frac{2\pi}{\ln2}
 \frac{\sqrt{2/(\gamma\varepsilon)}}{1-(2/r)\sqrt b}.
\end{equation}
At the boundary of the available packet budget, differentiation in
$b$ gives $b=(r/2)^{2/3}$. The resulting infimum within this
certificate family is
\[
 \frac{2\pi}{\ln2}
 \frac{\sqrt{8\log_2r/(\gamma(1-\delta))}}
 {(1-(2/r)^{2/3})^{3/2}}.
\]
The continuous stationary equation in $r$ is
$(r/2)^{2/3}=1+2\ln r$, near $r=54$. Our clean integers
$r=54$, $b=9$, $\varepsilon=1/260$ lie strictly inside the
allowed region for $\delta=9/10$. One also requires
$\varepsilon b/(b-1)<\delta$ so that new packet vertices respect
the degree cap; this holds for every numerical tuple used here.
This explains the choice without
claiming global optimality among contraction algorithms.

\section{Consequences for 3-SAT and density-sensitive inputs}\label{sec:sat}
For fixed clause density the direct clause/equality network has
$s=O_k(n+m)=O_{k,D}(n)$. For arbitrary density the source's
parity-query and incidence-index reduction has
$s=O_k(n\log\log n)$. Substitute these sizes in
Theorem~\ref{thm:main}; all surrounding overheads are lower order
in the corresponding exponential scale or part of the polynomial
factor. This improves the specified contraction certificate in both
reductions and preserves polynomial space.

A numerical coefficient of $133.207\ldots$ must not be written
directly in front of $n$ or $n\log\log n$ without tracking the
constant in the reduction's wire budget. The theorem is expressed
in the normalized network size precisely to avoid that ambiguity.

The existing scalar parity refinement can be kept at the reduction
stage. It removes a truncated polynomial coordinate and reduces a
retained incidence-degree bound, while our theorem permits $\ell=1$.
The two density packets can likewise keep their row-selection
procedures and their $O(n\log(2+B))$ wire budget for
$B\le\log^A n$ with fixed $A$. The new contraction then gives
the same density-sensitive order with the improved normalized
certificate. Neither composition creates a new arbitrary-density
log-power by itself.

\subsection{A further formula-specific degree choice}
The sparse-query proof groups tests and original variables separately
into $J=\ceil{n^{1/4}}$ groups. This exponent is not essential.
For example choose $J=\ceil{n^{0.495}}$. The total-wire estimate
remains $O(n\log(2+S))$ when each query has at most $S$
monomials and $S=\log^{O(1)}n$: the extra $J^2=n^{0.99+o(1)}$
rounding term is $o(n)$. Degree-array enumeration costs
$2^{O(n^{0.99}\log n)}$, which is lower order than
$2^{\Theta(n/\sqrt{\log n})}$. After the source's high-incidence
branching, each group has degree
$n^{0.505}\log^{O(1)}n\le s^{0.51}$ eventually.

The original entry-evaluation proof continues to apply. A test-group
entry is a product of independent per-test finite-difference
coefficients, each over a polylogarithmic scope. A variable-group
entry is the power of two determined by the number of untouched
variables. Increasing the number of groups does not create an
exponential product of entry runtimes~\cite{reduction}.

Consequently the same proof can use $\delta=0.51$ and
$\varepsilon=1/54$, with $r=54$ and $b=9$. Here
$\sum u_d\le\lambda/48$, and
\[
 0.51+\frac{\log_2 54}{12}<1
\]
supplies strict packing slack. Formula~\eqref{eq:parameter-family}
then gives the smaller network coefficient
\begin{equation}\label{eq:formula-specific}
 C_{\mathrm{query}}
 =\frac{9\pi}{4\ln2}\sqrt{\frac{108}{64/21}}
 =60.7069575490\ldots.
\end{equation}
The exact decimal is generated by the included scripts; the displayed
formula is authoritative. This secondary result requires the lower
degree promise and the stated regrouping. It is not the apples-to-apples
comparison against the generic $s^{9/10}$ interface in
Table~\ref{tab:constants}. Both have the same log-power.

\section{An additional exact variance lemma}\label{sec:variance}
The sharp-padding result already proves the main theorem using the
source's variance bound. A further structural refinement is available
when parallel-edge multiplicity is controlled.

\begin{lemma}[Exact collision variance]\label{lem:variance}
Give vertices on each side independent uniform labels from a set of
size $R$. Let $P$ be the number of ordered edge pairs with both endpoints
in common, including each edge paired with itself, and let $S$ count
ordered edge pairs sharing exactly one endpoint. If $X_\gamma$ is the
load of a pair of labels, then
\[
 \mathbb E\sum_\gamma(X_\gamma-M/R^2)^2
 =P(1-R^{-2})+S(R^{-1}-R^{-2}).
\]
In particular, if maximum degree is $\Delta$ and maximum parallel
multiplicity is $\mu$, the expectation is at most
$M\mu+2M\Delta/R$.
\end{lemma}
\begin{proof}
Expand the sum as the number of ordered edge pairs receiving identical
label pairs, minus $M^2/R^2$. For disjoint endpoints the agreement
probability is $R^{-2}$, for exactly one common endpoint it is
$R^{-1}$, and for two it is one. Subtracting $R^{-2}$ per pair
gives the equality. Also $P=\sum_{u,v}\mu_{uv}^2\le M\mu$
and $S\le2M\Delta$, which give the upper bound.
\end{proof}

With $R=r^a$, the corresponding relative padding error is bounded by
\[
 \sqrt{\frac{\mu R^4+2\Delta R^3}{M}}+\frac{LR^2}{M}.
\]
Thus $\mu R^4=o(M)$ and $\Delta R^3=o(M)$ can replace
the older $\Delta R^4=o(M)$ condition. A child has parallel
multiplicity at most $2L$ because one endpoint of every newly formed
edge is a packet vertex. A generic root, however, can have much
larger parallel multiplicity. We do not silently assume a simple root
or use this stronger variance estimate in the primary coefficient.
It is a separately proved improvement available to future instances
with an authenticated multiplicity bound.

\section{What remains unresolved}\label{sec:limits}
The new weight sequence is still linear:
$w_b=(64/21)b+O(\log b)$. Therefore its partition generating
function still has logarithm proportional to $1/\theta$.
Combining it with precision $P_t=\Theta(MD_i/u_d)$ still gives
a square-root saddle, namely $M/\sqrt{u_d}$. Packing and the
guard improve the constants and how the available logarithmic
budget is spent; they do not turn that saddle into $M/u_d$.

The old archive's binary-selection argument also remains relevant:
an abstract inventory with $\Theta(M/D)$ linear-weight modes has
$2^{\Omega(M/\sqrt u)}$ feasible tuples within an appropriate
budget. This is a barrier to improving the entire tuple-universe
count using only that numerical envelope, not a lower bound for
SAT or for non-enumerative contraction~\cite{analysispacket,cauchypacket}.
The new residue restriction changes the linear coefficient of the
weights and therefore escapes no such order-level barrier.

The dense-input $\log\log n$ has a separate origin. It encodes
degree-$O(n)$ subsets of an incidence universe of size
$n\,\polylog n$. Source and target group tensors remember which
test supplied an incidence, so replacing that code with the union
of original variables requires a new aggregation argument. The
previous degree-preserving compression, structured-mode summation,
and second-contraction proposals isolate this missing structure,
but none of the local improvements in this report constructs it.

For the distinct 49-page crossing-decomposition algorithm, a
stronger log-power would follow from admissible pieces with a
smaller contraction factor $\rho$, together with all uniform
description and evaluation bounds required in its Definition 7.
Its current partition-based construction has $\rho=1/\sqrt2$.
The present two-adic basis is not automatically a set of such
pieces. No unproved transfer between the two architectures is used.

\paragraph{Scope of novelty.}
The elementary finite Cauchy telescope is classical in character and
already appears in the factorial-Cauchy packet. The new use here is
the actual-row residue exclusion, the explicit $11{:}6{:}4$ pivot
allocation, and its compatible valuation gain. Sharp padding and
the additive guard were not present in the seven reviewed packets.
We establish novelty relative to that inspected corpus, not a
literature-wide priority claim.

\section{Executable artifacts and verification}\label{sec:checks}
The bundle contains the complete LaTeX source, its compiled PDF,
standard-library Python programs, recorded outputs, and a source/commit
manifest. The programs check finite exact statements underlying the new
lemmas. They are not a complete implementation of the preprint's SAT
counter. No SAT benchmark, complete Lean build, or transitive axiom
audit is claimed.

\begin{center}\small
\begin{tabularx}{\textwidth}{@{}X r@{}}
\toprule
Exact finite check & Cases\\
\midrule
Pivot-prefix counts, monotonicity and weight estimates & 10,000\\
Integer-domain pivot-product valuation inequalities & 8,505\\
Actual crossing-row parameter and residue cases & 19,008\\
\quad Included cases with non-prime-field interpolation & 9,504\\
Exact finite Cauchy separation identities & 50,496\\
Target-coefficient valuation inequalities & 239,808\\
Schedule guard and packet-budget identities & 25,704\\
Selected logarithmic sizes, with all provisional active depths & 9\\
Small multigraph/label-size configurations & 160\\
\quad Fully enumerated vertex labelings & 6,096\\
Finite abstract mode inventories & 21\\
\bottomrule
\end{tabularx}
\end{center}

The kernel, schedule and packing checks pass with ordinary Python and
with assertions disabled by \texttt{python3 -O}; their guards use explicit
exceptions. The result records preserve the scope and the code hashes.
The finite mode inventories use one, two or four colors and strict budgets
from 16 to 1024. In all 21 of these sampled inventories the old
factorial-Cauchy sequence has fewer admissible tuples than the new periodic
sequence. This is consistent with the finite-index qualification in
Section~\ref{sec:pivot21}: the theorem is asymptotic, and the periodic
sequence has lower-order costs. The optional exact-count selection keeps
the old basis whenever its finite inventory is smaller. These data are
included as a check against interpreting the asymptotic coefficient as a
uniform finite-input speedup.


The mode checks use exact rational arithmetic in
$\Q[U]/(U^2+U+1)$, including $\zeta=U$ and $1+U$. These
non-prime-field interpolation values matter: ordinary binomial
integrality is false for arbitrary elements of an unramified
extension. The proof and checks instead use the actual rows'
high-precision congruence to the prime subring.

To rerun the checks, execute \texttt{python3 code/run\_checks.py}
from the extracted archive. To rebuild the PDF, run
\texttt{pdflatex -no-shell-escape -halt-on-error report.tex} twice
from the archive root. Standard AMS, Latin Modern, geometry,
microtype, mathrsfs, needspace, hyperref, xurl, and table packages are needed. The
check suite itself requires only Python's standard library and
uses no network. The source manifest identifies the reviewed
external files without executing or bundling their programs.

\section{Research outcome}
The deliverable establishes an explicit improvement in the exponential
certificate of the reviewed OpenAI-style contraction, with fully stated
new local lemmas and preserved recursive interfaces. The clean generic
coefficient is $\newconstant$; under the explained lower-degree
sparse-query interface a further formula-specific coefficient follows.
The generic result already improves the strongest comparable previous
certificate, including its straightforward composition with factorial
Cauchy modes.

The stronger goal of improving
$2^{O(n\log\log n/\sqrt{\log n})}$ to a strictly smaller
exponent order remains open in this work. The exact point of progress
is therefore a proved quantitative refinement, not a completed new
log-power theorem.

\begin{thebibliography}{12}\small\raggedright
\bibitem{openai} OpenAI. \emph{The Exponential Time Hypothesis is False}.
Revision of 25 October 2026. Pinned repository commit \texttt{\opin}.
\href{https://github.com/openai/math/blob/7ef5c0cc5ca4b0e072b7a050fca09ab24bc6fa3a/preprints/The-Exponential-Time-Hypothesis-is-False-October-25-2026/paper.pdf}{Paper PDF}.
\bibitem{contraction} OpenAI. \emph{Proof of the contraction theorem},
\texttt{quantitative-contraction.tex}, same pinned commit.
\href{https://github.com/openai/math/blob/7ef5c0cc5ca4b0e072b7a050fca09ab24bc6fa3a/preprints/The-Exponential-Time-Hypothesis-is-False-October-25-2026/build/quantitative-contraction.tex}{Source}.
\bibitem{reduction} OpenAI. \emph{A sharper bound for exact counting},
\texttt{quantitative-reduction.tex}, same pinned commit.
\href{https://github.com/openai/math/blob/7ef5c0cc5ca4b0e072b7a050fca09ab24bc6fa3a/preprints/The-Exponential-Time-Hypothesis-is-False-October-25-2026/build/quantitative-reduction.tex}{Source}.
\bibitem{cauchypacket} ProveIt incoming research.
\emph{A factorial Cauchy expansion for the subexponential 3-SAT algorithm}.
\href{https://github.com/VladimirReshetnikov/ProveIt/blob/aa6933442eb936962de79ba6d321cd0391c450bd/docs/incoming/cauchy_mode_improvement_2026-10-26.zip}{Archived packet}.
\bibitem{schedulepacket} ProveIt incoming research.
\emph{Quantitative Improvements to the Exact-Counting Contraction Algorithm}.
\href{https://github.com/VladimirReshetnikov/ProveIt/blob/aa6933442eb936962de79ba6d321cd0391c450bd/docs/incoming/eth_contraction_improvements.zip}{Archived packet}.
\bibitem{densitypacket} ProveIt incoming research.
\emph{Density-aware exact SAT counting}.
\href{https://github.com/VladimirReshetnikov/ProveIt/blob/aa6933442eb936962de79ba6d321cd0391c450bd/docs/incoming/sat_density_improvement.zip}{Archived packet}.
\bibitem{densitypacket2} ProveIt incoming research.
\emph{ETH density refinement}.
\href{https://github.com/VladimirReshetnikov/ProveIt/blob/aa6933442eb936962de79ba6d321cd0391c450bd/docs/incoming/eth_density_refinement.zip}{Archived packet}.
\bibitem{precisionpacket} ProveIt incoming research.
\emph{Precision refinement for the 3-SAT contraction}.
\href{https://github.com/VladimirReshetnikov/ProveIt/blob/aa6933442eb936962de79ba6d321cd0391c450bd/docs/incoming/3sat_precision_refinement_2026-10-26.zip}{Archived packet}.
\bibitem{scalarpacket} ProveIt incoming research.
\emph{A Scalar Parity Refinement of the Subexponential SAT Counter}.
\href{https://github.com/VladimirReshetnikov/ProveIt/blob/aa6933442eb936962de79ba6d321cd0391c450bd/docs/incoming/3sat_refinement_package.zip}{Archived packet}.
\bibitem{analysispacket} ProveIt incoming research.
\emph{Asymptotic improvement analysis for the October 2026 3 SAT preprint}.
\href{https://github.com/VladimirReshetnikov/ProveIt/blob/aa6933442eb936962de79ba6d321cd0391c450bd/docs/incoming/3sat_asymptotic_analysis_20261027.zip}{Archived packet}.
\bibitem{otherpaper} Claude, Anthropic (attribution printed on manuscript).
\emph{A $2^{o(n)}$ time algorithm for 3-SAT}. 26 October 2026,
3:13 PM PT, 49 pages.
\href{https://github.com/VladimirReshetnikov/ProveIt/blob/aa6933442eb936962de79ba6d321cd0391c450bd/docs/incoming/eth.pdf}{Pinned supplied PDF}.
\bibitem{placement} ProveIt. Placement commit
\href{https://github.com/VladimirReshetnikov/ProveIt/commit/2b925d8800cda68f5723b65dc1fcdaf3d8170359}{\texttt{2b925d8800}}
and \href{https://github.com/VladimirReshetnikov/ProveIt/blob/2b925d8800cda68f5723b65dc1fcdaf3d8170359/Computability/SubexponentialSAT/README.md}{new project README},
26 October 2026, 5:58:36 PM Pacific time.
\end{thebibliography}
\end{document}
